\section{结果与讨论}

\subsection{Bi/MXene 复合材料的合成与结构表征}

Bi/MXene 复合材料通过两步路线合成，如示意图 1 所示。首先，以 Ti$_{3}$AlC$_{2}$ MAX 相为原料，采用 LiF/HCl 刻蚀剂在 40 $^{\circ}$C 下选择性刻蚀 Al 层 48 h，经冰浴超声处理后得到高浓度 MXene 胶体溶液（20--25 mg mL$^{-1}$）。随后，将 2 mmol BiCl$_{3}\cdot x$H$_{2}$O 溶于无水乙醇中，与冷冻干燥的 MXene 气凝胶（0.21 g）混合，磁力搅拌 3 h。将混合物转移至聚四氟乙烯内衬的不锈钢水热反应釜中，在 180 $^{\circ}$C 下反应 16 h，此过程中 Bi$^{3+}$ 发生水解，在 MXene 表面成核形成 Bi(OH)$_{3}$ 纳米颗粒。所得 Bi(OH)$_{3}$/MXene 中间产物经离心收集，用乙醇和超纯水洗涤后，在真空烘箱中 60 $^{\circ}$C 干燥。最后，将干燥粉末在管式炉中 Ar 气氛下 350 $^{\circ}$C 煅烧 2 h，Bi(OH)$_{3}$ 热分解并还原为金属 Bi，得到 Bi/MXene 复合材料。

通过扫描电子显微镜（SEM）对所制备材料的形貌进行了表征。图 Xa 展示了纯 MXene 纳米片的图像，显示出典型的层状形貌，表面具有明显的褶皱特征——为铋纳米颗粒的均匀负载提供了理想平台。在 Bi/MXene 复合材料中（图 Xb,c），可观察到大量的细小铋纳米颗粒均匀分散在 MXene 表面和层间通道内，形成了独特的"三明治"型结构。这种高度分散的纳米结构有望缩短 Na$^{+}$ 的扩散路径，同时 MXene 片层对铋颗粒的机械包覆可有效缓冲钠化/脱钠过程中的剧烈体积波动，从而增强电极结构的完整性。

能量色散 X 射线光谱（EDS）元素面分布图（图 X）为 Bi 在 MXene 骨架中的限域分布提供了直接证据。MXene 特征元素 Ti、O 和 F 与 Bi 信号在空间分布上表现出极强的相关性，证实了铋纳米颗粒被紧密整合到 MXene 层状基底中。这种紧密的界面接触构成了"三明治"限域结构的物理基础，使 MXene 能够在循环过程中对铋施加强有力的物理约束，从而大幅提升电极的结构稳定性和循环寿命。

通过高分辨率透射电子显微镜（HRTEM）进一步阐释了 Bi/MXene 复合材料的原子尺度界面结构。如图 X 所示，观察到铋纳米晶体被 MXene 纳米片紧密包裹。清晰的晶格条纹间距为 0.329 nm，对应于三方相铋的 (012) 晶面，与 XRD 结果一致。这种紧密的 Bi--MXene 界面接触为"三明治"限域架构提供了直接的形貌证据。在此构型中，MXene 片层发挥双重结构功能：既作为高导电基底加速电子传输，又通过力学性能强韧的二维骨架对铋颗粒产生强有力的物理限域作用。这一通过紧密 Bi--MXene 界面接触实现的纳米尺度限域效应，从根本上抑制了钠化/脱钠过程中的颗粒粉化，从而赋予电极卓越的结构稳定性和长循环耐久性。

通过 X 射线衍射（XRD）对 Bi/MXene 复合材料的晶体结构进行了表征。如图 X 所示，在 2$\theta$ = 27.2$^{\circ}$ 处出现强衍射峰，与三方相铋的标准卡片（PDF\#85-1329）的 (012) 晶面完全吻合，证实了铋以结晶良好的金属态存在。与此同时，在 2$\theta \approx$ 8$^{\circ}$ 处观察到 MXene 的特征 (002) 晶面衍射峰，表明在水热和煅烧处理后，Ti$_{3}$C$_{2}$T$_{x}$ 的层状结构得以保持。

\subsection{电化学储钠性能}

采用扣式半电池评估了 Bi/MXene 复合材料的电化学储钠行为。图 X 展示了不同循环圈数下的恒流充放电（GCD）电压曲线。在放电过程中，可观察到对应于铋与钠合金化反应生成 Na$_{3}$Bi 的明确电压平台。值得注意的是，从第 3 次到第 100 次循环的充放电曲线几乎完全重合，表明 Bi/MXene 电极具有极高的电化学可逆性。此外，充电和放电支路之间极小的电压迟滞为 Bi/MXene 界面处快速的离子传输动力学和较低的电荷转移电阻提供了直接证据。

通过测量 0.1 至 10 A g$^{-1}$ 电流密度范围内的比容量（图 X），评估了 Bi/MXene 电极的倍率性能。在 0.1、0.2、0.5、1、2、5 和 10 A g$^{-1}$ 的电流密度下，电极的可逆容量分别为 356.6、354.8、351.6、348.9、346.2、341.8 和 336.9 mAh g$^{-1}$。值得注意的是，将电流密度提升 100 倍（从 0.1 至 10 A g$^{-1}$）仅导致 5.5\% 的容量下降，充分体现了 Bi/MXene 复合材料卓越的倍率耐受能力。更重要的是，当电流密度恢复至 0.1 A g$^{-1}$ 时，比容量恢复至 362.9 mAh g$^{-1}$——略高于初始值 356.6 mAh g$^{-1}$，这表明电极在倍率测试过程中经历了逐步活化。经历高倍率冲击后容量的完全恢复，证明了 Bi/MXene 电极优异的结构稳健性。这一卓越的倍率性能归因于 MXene 骨架提供的高效电子传输网络以及铋纳米颗粒大幅缩短的 Na$^{+}$ 扩散距离。

在 2 A g$^{-1}$ 的电流密度下评估了 Bi/MXene 复合材料的长循环稳定性（图 X）。首次循环中，电极的放电容量为 422.9 mAh g$^{-1}$，充电容量为 372.6 mAh g$^{-1}$，初始库仑效率为 88.1\%。首次循环中约 50 mAh g$^{-1}$ 的不可逆容量损失归因于固体电解质界面膜（SEI）的形成和 Na$^{+}$ 在 MXene 残余表面官能团中的俘获。经过前三圈的初始活化阶段后，电极稳定在 343.0 mAh g$^{-1}$ 的可逆容量（第 4 圈）。在长期循环中，电极展现出卓越的耐久性，历经 5,000 次连续充放电后仍保持约 310 mAh g$^{-1}$ 的比容量，相对于第 4 圈稳定容量的保持率为 90.4\%。库仑效率在初始循环后迅速提升，并在整个后续循环期内稳定在接近 100\%。这一出色的循环稳定性归功于 MXene 骨架的"三明治"限域效应，有效缓冲了铋在钠化/脱钠过程中的剧烈体积变化，抑制了活性材料的粉化与脱落，从而保证了电极结构的长期完整性。值得指出的是，Bi/MXene 复合材料的循环稳定性与近期报道的高性能铋-碳负极相比毫不逊色——后者虽能在数千次循环中达到相当的保持率，但通常需依赖氮掺杂碳封装或多组分异质结构设计 [10,11,17]。

\subsection{反应动力学与电荷存储机理}

为阐明卓越倍率性能背后的电荷存储机理，在不同扫速（0.2 至 5.0 mV s$^{-1}$）下进行了循环伏安（CV）测试（图 X）。所有扫速下清晰的氧化还原峰证实了铋与钠之间高度可逆的合金化/去合金化反应。根据幂律方程 $i = av^{b}$ 分析了峰电流 ($i$) 与扫速 ($v$) 的关系，其中 $b$ 值可区分扩散控制的法拉第过程（$b \approx$ 0.5）和表面限域的电容过程（$b \approx$ 1.0）。对数拟合 $\log i$ 对 $\log v$ 得到 $b$ 值在 0.62--0.82 范围内，表明电荷存储通过扩散控制的电池型行为和表面限域赝电容贡献共同参与的混合机制进行。对电流响应的定量解耦分析表明，在所测扫速范围内，扩散控制过程占主导地位，贡献了总容量的 92\%--98\%。尽管如此，由 MXene 骨架提供的不可忽略的电容贡献在维持高电流密度下的快速电荷转移中发挥了关键作用，从而使电极在极端倍率条件下仍能保持高容量。

采用电化学阻抗谱（EIS）对比研究了 Bi/MXene 复合材料与纯铋电极的电荷传输特性（图 X）。Bi/MXene 电极的 Nyquist 图在中高频区域呈现出远小于纯铋电极的半圆，表明其电荷转移电阻（$R_{\text{ct}}$ = 0.69 $\Omega$）显著低于纯铋电极（$R_{\text{ct}}$ = 13.95 $\Omega$）。此外，Bi/MXene 电极在低频区域展现出更为陡峭的 Warburg 斜率，表明电极内部 Na$^{+}$ 扩散更为顺畅。这些结果有力地证明了二维 MXene 导电骨架与铋纳米颗粒的紧密接触，构建了高效的电子传输通道，显著降低了界面电荷转移势垒，加速了电极反应动力学。

采用恒电流间歇滴定技术（GITT）定量评估了整个钠化/脱钠电压区间内的 Na$^{+}$ 扩散系数（$D_{\text{Na}^{+}}$）（图 X）。Bi/MXene 电极在整个电压窗口内表现出较高的 $D_{\text{Na}^{+}}$，数值范围在 $10^{-12.5}$--$10^{-13}$ cm$^{2}$ s$^{-1}$ 之间。值得注意的是，GITT 测试中的极化电压仅为 74.6 mV，反映出极低的欧姆过电位和浓度过电位。这一结果直接验证了"三明治"限域结构的设计原理：MXene 骨架增强了电子导电性，而纳米化的铋颗粒大幅缩短了固态 Na$^{+}$ 扩散距离，从而在原子尺度上实现了快速的离子/电子协同传输。

\subsection{循环后结构分析}

为评估 Bi/MXene 电极在长时间电化学循环后的结构完整性，对在 2 A g$^{-1}$ 下循环 5,000 次后的电极进行了非原位 SEM 表征（图 X）。与循环后铋负极常见到的严重粉化和团聚现象形成鲜明对比 [9,10]，Bi/MXene 电极在很大程度上保留了原始形貌，铋纳米颗粒仍良好分散于 MXene 层间通道内，未见大面积的裂纹或剥离。这种形态学稳定性从视觉上直接证实了 MXene 的"三明治"限域效应有效抑制了传统合金型负极普遍存在的机械降解途径——颗粒粉化、电接触失效及 SEI 反复重构。以上循环后观察结果与动力学分析共同确立：Bi/MXene "三明治"结构在维持快速离子/电子传输的同时，实现了可靠的电极级结构稳定性，从而为高倍率性能和长循环耐久性的同步达成奠定了基础。

\section{结论}

综上所述，本文报道了一种具有独特"三明治"结构的二元 Bi/MXene 复合负极，通过简便的水热路线结合热还原法合成。SEM、HRTEM、EDS 元素面分析和 XRD 表征证实，铋纳米颗粒被均匀限域在 Ti$_{3}$C$_{2}$T$_{x}$ MXene 纳米片之间。这一结构简单的设计无需硫化、碳包覆或多组分异质结工程，即可同时实现卓越的倍率性能和超长循环稳定性。

Bi/MXene 电极在 0.1 A g$^{-1}$ 下实现 356.6 mAh g$^{-1}$ 的可逆容量，在 10 A g$^{-1}$ 下仍保持 336.9 mAh g$^{-1}$，对应于电流密度提升 100 倍后 94.5\% 的容量保持率。在长期循环中，电极在 2 A g$^{-1}$ 下经 5,000 次循环后维持约 310 mAh g$^{-1}$ 的容量，容量保持率达 90.4\%，库仑效率趋于 100\%。初始库仑效率为 88.1\%。

机理研究表明，卓越的电化学性能源自结构优势与动力学优势的协同组合。MXene 骨架提供连续的类金属导电网络，将电荷转移电阻从纯铋的 13.95 $\Omega$ 降至 0.69 $\Omega$；同时，其力学性能强韧的二维骨架对铋纳米颗粒实施物理限域，抑制了传统合金型负极中导致容量衰减的粉化和电接触失效。CV 动力学分析确认了混合电荷存储机理（$b$ = 0.62--0.82），其中 MXene 带来的赝电容贡献在补充铋的扩散控制容量的同时，维持了高倍率下的快速电荷转移。GITT 测试进一步印证了上述发现：Na$^{+}$ 扩散系数在 $10^{-12.5}$--$10^{-13}$ cm$^{2}$ s$^{-1}$ 范围内，且极化电压仅为 74.6 mV。循环后 SEM 则直接验证了"三明治"结构在 5,000 次循环后仍能保持电极形貌完好，与本文强调的限域驱动结构稳定性的核心设计理念相吻合。

这些结果表明，合金型钠离子电池负极长期面临的倍率-容量-稳定性权衡困境可通过结构简单、可规模化制备的二元复合结构加以克服——为高性能快充钠离子电池开辟了一条切实可行的技术路径。
